\section{Introduction}
\label{sec:introduction}

A student working through a lecture may need an explanation of a formula,
a small example, and an opportunity to apply the idea without help.
These activities also create information that could guide the next study
session. For an instructor, the useful question is more specific than
whether students completed a unit: which concept caused difficulty, and
which part of the course material should be revisited?

StudyBirds is designed around this connection between course content,
student activity, and teaching feedback. A course provides the shared
context. Students choose a topic and a study purpose, work with a tutor,
and retain their attempts and session outputs. The instructor-facing
design uses permitted evidence to identify difficulties associated with
particular concepts and source locations. This creates a common reference
point for learning support and teaching decisions.

The research question is whether this connected workflow can support
useful, inspectable guidance while preserving students' independent
reasoning. A correct answer after a hint, for example, should remain
distinguishable from a correct independent attempt. Similarly, an
instructor insight should communicate what evidence supports it instead
of treating a generated explanation as an established diagnosis.

This draft develops three starting points for the paper: a course-based
workflow joining preparation and study; an explicit distinction between
assisted and independent performance; and a proposed evaluation connecting
system behavior to student and instructor tasks. These are design
contributions to be examined, rather than established claims of novelty
or improved learning.
